\chapter{A Concrete Naming Certificate for $\DyadicDivisionUp$}
\label{ch:concrete-instances-dyadic-division-namecert}
\origin{human}

Following Bishop and Bridges constructive dyadic arithmetic, $\DyadicDivisionUp$ names the finite quotient surface used when an exposed dyadic numerator is divided by a displayed nonzero dyadic denominator before regular rational and real-name consumers can read the result. It sits over the dyadic arithmetic carrier of \autoref{ch:concrete-instances-dyadicratcore-namecert}, the comparison and order rows of \autoref{ch:concrete-instances-dyadic-comparison-namecert} and \autoref{ch:concrete-instances-dyadic-order-namecert}, the reciprocal-window budget of \autoref{ch:concrete-instances-cauchyinversebudget-namecert}, the regular rational handoff of \autoref{ch:concrete-instances-regseqrat-namecert}, and the terminal real-name seal of \autoref{ch:concrete-instances-real-namecert}. The packet records division as displayed finite NameCert data; it is not quotient rational equality, a total inverse operation, trichotomy, selected representatives, countable choice, or ambient $\RealUp$ completeness.

\begin{definition}[Dyadic division carrier]
\label{def:dyadic-division-carrier}
A $\DyadicDivisionUp$ carrier is a finite $\BHist$ packet
$$
Q=(A,B,L,R,T,E,H,C,P,N)
$$
with the following displayed rows:
\begin{enumerate}
\item $A$ is the numerator dyadic packet, read through the $\DyadicRatCoreUp$ arithmetic surface with its displayed sign, scale, endpoint, and tolerance data.
\item $B$ is the nonzero denominator dyadic packet, also read through $\DyadicRatCoreUp$, whose denominator role is fixed before any reciprocal or quotient row is consumed.
\item $L$ is the denominator apartness and bound row. It records the finite comparison and order evidence, supplied by $\DyadicComparisonUp$ or $\DyadicOrderUp$, that keeps $B$ away from the zero boundary and exposes a usable lower bound.
\item $R$ is the reciprocal ledger. It is the $\CauchyInverseBudgetUp$-compatible finite row that converts the displayed nonzero denominator packet and apartness bound into reciprocal-read data without exporting a total inverse operation.
\item $T$ is the quotient readback row. It reads the numerator row $A$ and reciprocal ledger $R$ as one displayed dyadic quotient packet before any regular rational handoff is available.
\item $E$ is the $\RegSeqRatUp$ handoff together with the terminal $\RealUp$ seal row. It records the regular rational readback of $T$ and the downstream real-name seal only after $A,B,L,R,T$ have been displayed.
\item $H$ records componentwise $\hsame$ transport for numerator, denominator, apartness, reciprocal, quotient, readback, seal, replay, provenance, and local naming rows.
\item $C$ records displayed $\Cont$ replay from numerator and denominator inspection through apartness exposure, reciprocal ledger construction, quotient readback, $\RegSeqRatUp$ handoff, $\RealUp$ sealing, transport, provenance, and local naming.
\item $P$ is the $\Pkg$ provenance over $\DyadicDivisionUp$, $\DyadicRatCoreUp$, $\DyadicComparisonUp$, $\DyadicOrderUp$, $\CauchyInverseBudgetUp$, $\RegSeqRatUp$, $\RealUp$, $\BHist$, $\hsame$, $\Cont$, and $\NameCert$.
\item $N$ is the local $\NameCert_{\DyadicDivisionUp}$ row.
\end{enumerate}
The classifier compares accepted packets only through the tuple $A,B,L,R,T,E,H,C,P,N$. It has no coordinate for quotient rational equality, a total inverse operation, trichotomy, selected representatives, countable choice, ambient $\RealUp$ completeness, Lean verification, or bridge checking.
\end{definition}

\begin{theorem}[Dyadic division NameCert obligations]
\label{thm:dyadic-division-namecert-obligations}
For every accepted $\DyadicDivisionUp$ carrier $Q=(A,B,L,R,T,E,H,C,P,N)$, the local $\NameCert_{\DyadicDivisionUp}$ surface consists exactly of the following finite obligation rows:
\begin{enumerate}
\item carrier admission for the displayed numerator packet, nonzero denominator packet, denominator apartness and bound row, reciprocal ledger, quotient readback, regular rational handoff, real seal, transport, replay, provenance, and local naming rows;
\item dyadic-source exposure, requiring $A$ and $B$ to be read as $\DyadicRatCoreUp$ packets before any quotient row is available;
\item denominator guard exposure, requiring $L$ to display comparison or order evidence keeping $B$ away from zero before $R$ is consumed;
\item reciprocal-ledger locality, requiring $R$ to use only the displayed denominator packet, apartness bound, and reciprocal-window budget named by $\CauchyInverseBudgetUp$;
\item quotient readback exactness, requiring $T$ to consume $A$ and $R$ together and to expose no independent quotient representative;
\item $\RegSeqRatUp$ and $\RealUp$ non-escape, requiring $E$ to remain downstream of $A,B,L,R,T$;
\item structural replay, requiring $H,C,P,N$ to preserve, replay, source, and name exactly the same finite packet.
\end{enumerate}
No dyadic division obligation is stored outside these rows.
\end{theorem}
\begin{proof}
Project the carrier of \autoref{def:dyadic-division-carrier}. The analytic rows are $A,B,L,R,T,E$, and the structural rows $H,C,P,N$ transport, replay, source, and name the same tuple.

The source exposure clauses are read from the sibling dyadic surfaces. The numerator and denominator rows are $\DyadicRatCoreUp$ packets, while the guard row $L$ is the finite comparison or order evidence supplied by \autoref{ch:concrete-instances-dyadic-comparison-namecert} or \autoref{ch:concrete-instances-dyadic-order-namecert}.

The reciprocal clause is then local. Since $R$ is downstream of $B$ and $L$, it has the same guarded shape as the reciprocal budget of \autoref{ch:concrete-instances-cauchyinversebudget-namecert}; no route to $R$ starts from an unguarded denominator.

The quotient row $T$ consumes $A$ and $R$ together. Therefore the quotient read is a displayed dyadic row rather than a selected quotient representative or a host rational equality class.

Finally $E$ is downstream of the quotient readback, and $H,C,P,N$ add only transport, replay, provenance, and local naming. Because none of the excluded principles is a carrier coordinate, no local NameCert obligation can be stored outside $A,B,L,R,T,E,H,C,P,N$.
\end{proof}

\begin{theorem}[Dyadic division nonzero reciprocal route]
\label{thm:dyadic-division-nonzero-reciprocal-route}
Let $Q=(A,B,L,R,T,E,H,C,P,N)$ be an accepted $\DyadicDivisionUp$ carrier. Every reciprocal ledger read used by the quotient row factors through
$$
B \longmapsto L \longmapsto R .
$$
Thus a denominator packet reaches reciprocal data only after the displayed apartness and bound row has been exposed, and a route omitting $L$ cannot export a quotient readback.
\end{theorem}
\begin{proof}
Begin with the denominator row $B$. Carrier admission places it before the apartness and reciprocal rows, so any denominator-facing reciprocal route starts from a displayed $\DyadicRatCoreUp$ packet.

Next read the guard row $L$. It is the only coordinate that records denominator apartness and a usable lower bound, and it is the only place where the $\DyadicComparisonUp$ or $\DyadicOrderUp$ evidence enters the packet.

Now project the reciprocal ledger $R$. Its construction is downstream of $B$ and $L$ and follows the reciprocal-window budget discipline named by \autoref{ch:concrete-instances-cauchyinversebudget-namecert}.

If $L$ is omitted, the carrier has no row certifying that the denominator stays away from zero. The route therefore remains at the boundary surface and cannot reach the quotient readback row $T$.
\end{proof}

\begin{theorem}[Dyadic division Real readback nonescape]
\label{thm:dyadic-division-real-readback-nonescape}
Every $\RegSeqRatUp$, $\RealUp$, L10 Real, or Completion consumer of an accepted $\DyadicDivisionUp$ carrier reads the displayed chain
$$
A,B,L \longmapsto R \longmapsto T \longmapsto E \longmapsto H,C,P,N .
$$
Consequently the consumer receives a finite dyadic quotient readback before the regular rational handoff or real seal is available, with no quotient rational equality, total inverse operation, trichotomy, selected representative, countable-choice row, or ambient $\RealUp$ completeness witness inserted into the route.
\end{theorem}
\begin{proof}
Project the quotient row $T$. It can be consumed only after the numerator packet $A$ and the guarded reciprocal ledger $R$ have been displayed.

Use \autoref{thm:dyadic-division-nonzero-reciprocal-route}. The reciprocal ledger itself has already passed through the denominator packet $B$ and the apartness row $L$, so the quotient route contains the complete nonzero-denominator guard.

The regular rational and real-facing row $E$ is terminal relative to $T$. It hands the displayed quotient row to $\RegSeqRatUp$ and then to the $\RealUp$ seal; it has no second entrance from a host quotient or total inverse operation.

The remaining rows $H,C,P,N$ transport, replay, source, and name the same chain. Since the carrier contains no excluded quotient, choice, trichotomy, or completeness coordinate, every L10 Real or Completion consumer remains inside the displayed route.
\end{proof}

\closureat{DyadicDivisionUp}{seedStr}

\begin{closurestatus}{\DyadicDivisionUp}
  \constructivestory{A $\DyadicDivisionUp$ packet is generated from a numerator DyadicRatCore packet, a nonzero denominator DyadicRatCore packet, a denominator apartness and bound row, a reciprocal ledger, a quotient readback row, a RegSeqRat handoff with terminal RealUp seal, componentwise hsame transport, displayed Cont replay, Pkg provenance, and a local NameCert row.}
  \theoryclosure{\seedClosure}
  \scopeclosed{\autoref{def:dyadic-division-carrier}, \autoref{thm:dyadic-division-namecert-obligations}, \autoref{thm:dyadic-division-nonzero-reciprocal-route}, and \autoref{thm:dyadic-division-real-readback-nonescape} seed the finite dyadic quotient surface for L10 Real and Completion consumers over $\DyadicRatCoreUp$, $\DyadicComparisonUp$, $\DyadicOrderUp$, $\CauchyInverseBudgetUp$, $\RegSeqRatUp$, and $\RealUp$.}
  \formalstatus{\unformalizedV}
  \bridgestatus{none}
  \notclaimed{This chapter does not close quotient rational equality, total inverse operations, trichotomy, selected representatives, countable choice, ambient $\RealUp$ completeness, Lean verification, or bridge checking.}
  \upgradepath{Obligation closure requires checked carrier admission, dyadic-source exposure, denominator guard exposure, reciprocal-ledger locality, quotient readback exactness, RegSeqRat and Real non-escape, structural replay, provenance, and local naming for BEDC.Derived.DyadicDivisionUp.}
\end{closurestatus}
